\sectioncentered*{Перечень условных обозначений}
\addcontentsline{toc}{section}{Перечень условных обозначений}
\label{sec:reduction}

\noindent
\begin{tabular}{@{}p{3.6cm}@{\ \mdash\ }p{11.0cm}@{}}
NIKA & Intelligent Knowledge-driven Assistant, интеллектуальный ассистент на базе OSTIS;\\
БЗ & База знаний;\\
SCs & Symbolic Code compact, текстовый язык представления SC-кода;\\
агент & Программный модуль, инициируемый событием в SC-памяти;\\
действие & Узел класса \texttt{actions\_class}, определяющий работу агента;\\
намерение & Класс сообщения, определяемый классификатором (intent);\\
сущность & Элемент БЗ, выделенный из сообщения пользователем (entity);\\
заметка & Отношение <<заметка>> (\texttt{nrel\_note}), текстовое описание сущности;\\
основной идентификатор & Отношение \texttt{nrel\_main\_idtf}, имя сущности на языке;\\
MediaWiki API & HTTP-интерфейс Википедии, отвечает в формате JSON;\\
JSON & JavaScript Object Notation, формат передачи данных;\\
witai-shim & Модуль классификации сообщений системы NIKA.
\end{tabular}
